\section{Группы}

\subsection{Введение}

\begin{defn}{Множество с бинарной операцией}{}
    Множество \( M \), с заданным отображением \( M \times M \to M \) \( (a, b) \mapsto a \circ b \)

    Обозначается \( (M, \circ) \)
\end{defn}

\begin{defn}{Полугруппа}{}
    Множество с бинарной операцией \( (M, \circ) \) называется полугруппой,
    если эта операция ассоциативна:
    \[
        \forall a, b, c \in M \ (a \circ b) \circ c = a \circ (b \circ c)
    \]

    Обозначается \( (S, \circ) \)
\end{defn}

\begin{example}{Не всякая операция ассоциативна}{}
    \begin{itemize}
        \item
            \( M = \NN \) и \( a \circ b = a^b \) --- не ассоциативна:
            \[
                2^{(1^2)} \neq (2^1)^2
            \]
        \item
            \( M = \ZZ \) и \( a \circ b = a - b \) --- тоже не ассоциативна
    \end{itemize}
\end{example}

\begin{defn}{Моноид}{}
    Полугруппа \( (S, \circ) \) называется моноидом,
    если в ней существует нейтральный элемент:
    \[
        \exists e \in S: \ e \circ a = a \circ e = a
    \]
\end{defn}

\begin{example}{}{}
    Является ли \( (\NN, +) \) моноидом?
    Ответ зависит от географического положения,
    ведь в России не \( 0 \) не считают натуральным,
    а, например, во Франции, наоборот.
    Далее в курсе \( 0 \) не является натуральным.
\end{example}

\begin{remark}{}{}
    Если нейтральный элемент существует, то он единственный:
    \[
        e_1 = e_1 \circ e_2 = e_2
    \]
\end{remark}

\begin{defn}{Группа}{}
    Моноид \( (S, \circ) \) называется группой,
    если \( \forall a \in S \: \exists b \in S : a \circ b = b \circ a = e \)

    Обозначается \( (G, \circ) = (G, \cdot) = G \).

    Чаще всего используется мультипликативная запись: \\
    \( (a, b) \mapsto ab\), \( e = 1 \), обратный элемент обозначается \( a^{-1} \).
\end{defn}

\begin{exercise}{}{}
    Если обратный элемент существует, то он единственный.
\end{exercise}

\begin{defn}{Абелева группа}{}
    Группа \( G \) коммутативна (абелева), если \( \forall a, b \in G \ ab = ba \).

    В абелевых группах используется аддитивное обозначение: \\
    \( (a, b) \mapsto a + b\), \( e = 0 \), обратный элемент обозначается \( -a \).
\end{defn}

\begin{defn}{Порядок группы}{}
    Порядок группы \( G \) \( |G| \) --- число элементов в ней.

    \begin{itemize}
        \item
            \( |G| < \infty \Rightarrow G \) --- конечна
        \item
            \( |G| = \infty \Rightarrow G \) --- бесконечна
    \end{itemize}
\end{defn}

\begin{example}{Группы}{}
    \begin{enumerate}
        \item
            Числовые аддитивные:

            \( (\ZZ, +) \), \( (\QQ, +) \), \( (\RR, +) \), \( (\CC, +) \), \( (\ZZ_n, +) \)
        \item
            Числовые мультипликативные:

            \( (\QQ \setminus \{ 0 \}, \times) \),
            \( (\RR \setminus \{ 0 \}, \times) \),
            \( (\CC \setminus \{ 0 \}, \times) \),
            \( (\ZZ_p \setminus \{ 0 \}, \times) \)
        \item
            Группы матриц:

            \( GL_n (\RR) = \{ A \in M_n(\RR) \ \vline \ \det A \neq 0 \} \) --- общая линейная

            \( SL_n (\RR) = \{ A \in M_n(\RR) \ \vline \ \det A = 1 \} \) --- специальная линейная

            Есть еще куча разных групп: верхнетреугольные / нижнетреугольные / диагональные
            (конечно же, нужно взять только обратимые)
        \item
            Группы перестановок:

            \( S_n \), \( |S_n| = n! \)

            \( A_n \) --- четные перестановки, \( |A_n| = \frac{n!}{2} \)
    \end{enumerate}
\end{example}

\begin{exercise}{}{}
    \begin{itemize}
        \item
            \( S_n \) коммутативна \( \Leftrightarrow n \leq 2 \)
        \item
            \( A_n \) коммутативна \( \Leftrightarrow n \leq 3 \)
    \end{itemize}
\end{exercise}

\begin{defn}{Подгруппа}{}
    Подмножество \( H \) группы \( G \) называется подгруппой,
    если (определения эквивалентны):
    \begin{itemize}
        \item
            \begin{enumerate}
                \item
                    \( H \neq \varnothing \)
                \item
                    \( \forall a, b \in H \: ab^{-1} \in H \).
            \end{enumerate}
        \item
            \begin{enumerate}
                \item
                    \( e \in H\)
                \item
                    \( \forall a, b \in H \: ab \in H \)
                \item
                    \( \forall a \in H \: a^{-1} \in H \)
            \end{enumerate}
    \end{itemize}
\end{defn}

\begin{remark}{}{}
    В любой группе есть две подгруппы \( H = \{ e \} \) и \( H = G \)

    Они называются несобственными.
\end{remark}

\begin{prop}{}{}
    Любая подгруппа в \( (\ZZ, +) \) имеет вид \( k \ZZ \) (\( k \in \ZZ_{\geq 0} \))
\end{prop}

Очевидно, что \( k \ZZ \) --- подгруппа.

Пусть \( H \subseteq \ZZ \) --- подгруппа.

Если \( H = \{ 0 \} \), то \( k = 0 \).

Иначе в \( H \) есть ненулевой элемент.
Тогда в \( H \) есть положительный элемент.
Пусть \( k \) --- наименьший положительный элемент в \( H \).
Тогда \( k \ZZ \subseteq H \).
Пусть \( a \in H \).
Поделим с остатком: \( a = kq + r \), где \( 0 \leq r < k \).
\( r = a - kq \). \( a \in H, k \in H \Rightarrow r \in H \).
\( k \) наименьшее положительное число в \( H \), значит \( r = 0 \),
то есть \( a \in k \ZZ \).

\begin{defn}{Циклическая подгруппа}{}
    Пусть \( G \) --- группа, \( g \in G \).

    Циклической подгруппой,
    порожденной элементом \( g \) называется
    \[
        H = \langle g \rangle = \{ g^m, m \in \ZZ \}
    \]
\end{defn}

\begin{defn}{Порядок элемента}{}
    Пусть \( G \) --- группа, \( g \in G \)

    Порядком элемента \( g \) называется наименьшее натуральное число \( m \),
    для которого \( g^m = e \) или бесконечность, если такого \( m \) не существует.

    Обознается \( \ord (g) \)
\end{defn}

\begin{prop}{}{}
    \[
        | \langle g \rangle | = \ord (g)
    \]
\end{prop}

Пусть \( g^n = g^m \) и \( n > m \Rightarrow g^{n - m} = e \Rightarrow \) порядок \( g \) конечен.

Отсюда можно сделать вывод, что если \( \ord (g) = \infty \), то \( g^n \neq g^m \) при \( n \neq m \),
а значит и \( | \langle g \rangle | = \infty \).

Теперь пусть \( \ord (g) = m \).
По определению \( g^m = e \) и \( g^k \neq e \) при \( 0 < k < m \).
Поэтому \( g^0, \ldots, g^{m - 1} \) попарно различны (иначе можно применить трюк выше).
Любое целое число можно поделить с остатком на \( m \): \( s = mq + r \), где \( 0 \leq r < m \).
Но тогда понятно, что \( g^s = g^r \).
А значит \( \langle g \rangle = \{ g^0, g^1, \ldots, g^{m - 1} \} \Rightarrow |G| = m \).

Что и требовалось доказать.

\begin{defn}{Циклическая группа}{}
    Группа \( G \) назыается циклической, если \( \exists g \in G : G = \langle G \rangle \)
\end{defn}

\begin{remark}{}{}
    Нетрудно убедиться, что тогда \( G \) коммутативна и не более чем счетна.
\end{remark}

\begin{example}{}{}
    \begin{itemize}
        \item
            \( (\ZZ, +) = \langle 1 \rangle = \langle -1 \rangle \)
        \item
            \( (\ZZ_n, +) = \langle \ol{1} \rangle = \langle \ol{k} \rangle \), где \( (k, n) = 1 \)
        \item
            \( (\CC \setminus \{ 0 \}, \times) \supset \langle \varepsilon \rangle \),
            где \( \varepsilon \) --- первообразный корень
    \end{itemize}
\end{example}

\begin{defn}{Левый смежный класс}{}
    Пусть \( G \) --- группа, \( H \subseteq G \) --- подгруппа.

    Левый смежный класс элемента \( g \in G \) по подгруппе \( H \) --- это подмножество,
    которое обозначается \( gH \) и равно \( \{ gh \: \vline \: h \in H \} \).
\end{defn}

\newpage

\begin{lemma}{}{}
    Либо \( g_1 H = g_2 H \), либо \( g_1 H \cap g_2 H = \varnothing \).
\end{lemma}

Если их пересечение пусто, то условие верно, иначе есть общий элемент:
\[
    g_1 h_1 = g_2 h_2 \Rightarrow g_1 = g_2 h_2 h_1^{-1}
\]

\( H \) --- подгруппа, поэтому \( h_2 h_1^{-1} \in H \Rightarrow g_1 \in g_2 H \Rightarrow g_1 H \subseteq g_2 H \).

Аналогично \( g_2 H \subseteq g_1 H \), откуда \( g_1 H = g_2 H \).

\begin{lemma}{}{}
    \[
        |gH| = |H|
    \]
\end{lemma}

\( gH = \{ gh \: \vline \: h \in H \} \),
поэтому давайте сопоставим \( h \in H \) элемент \( gh \).
Откуда \( |H| \geq |gH| \).

Осталось показать, что никакие два элемента в \( gH \) не совпадают:
\[
    g h_1 = g h_2 \Rightarrow h_1 = h_2
\]

А значит \( |H| = |gH| \), что и требовалось доказать.

\begin{defn}{Индекс подгруппы}{}
    Индекс подгруппы \( H \) в группе \( G \) --- это
    \[
        [ G : H ] = | \fc{G}{H} | = \text{число левых смежных классов}
    \]
\end{defn}

\begin{thrm}{}{}
    Если \( G \) --- конечная группа, а \( H \) в ней подгруппа,
    то \( |G| = |H| \cdot [G : H] \)
\end{thrm}

Явно следует из двух предыдущих лемм.

\begin{coroll}{}{}
    Если \( G \) конечно, и \( H \subseteq G \) --- подгруппа,
    то \( |G| \vdots |H| \).
\end{coroll}

\begin{coroll}{}{}
    Если \( G \) конечно, и \( H \subseteq G \) --- подгруппа,
    то \( \forall g \in G \: |G| \vdots \ord(g) \).
\end{coroll}

Явно следует из \( \ord (g) = | \langle g \rangle | \)

\begin{coroll}{}{}
    \[
        g^{|G|} = e
    \]
\end{coroll}

Явно следует из того, что \( |G| \vdots \ord (g) \)

\begin{coroll}{}{}
    Малая теорема Ферма
\end{coroll}

\begin{coroll}{}{}
    Если \( |G| = p \) --- простое, то \( G \) --- циклическая группа,
    порожденная своим любым неединичным элементом.
\end{coroll}

Пусть \( g \in G \) и \( g \neq e \).

Тогда \( p = |G| \vdots | \langle g \rangle | \).
Так как \( | \langle g \rangle | \geq 2 \), \( | \langle g \rangle | = p = |G| \),
а значит \( |G| = \langle g \rangle \).

\begin{defn}{Правый смежный класс}{}
    Пусть \( G \) --- группа, \( H \subseteq G \) --- подгруппа.

    Правый смежный класс элемента \( g \in G \) по подгруппе \( H \) --- это подмножество,
    которое обозначается \( Hg \) и равно \( \{ hg \: \vline \: h \in H \} \).
\end{defn}

\begin{remark}{}{}
    Для них верны те же свойства, что и для левых смежных классов.
\end{remark}

\begin{defn}{Нормальная подгруппа}{}
    Подгруппа \( H \) нормальна, если \( gH = Hg \) \( \forall g \in G \).
\end{defn}

\begin{prop}{}{}
    Для подгруппы \( H \subseteq G \) следующие условия эквивалентны:
    \begin{enumerate}
        \item
            \( H \) нормальна
        \item
            \( gHg^{-1} \subseteq H \ \forall g \in G \)
        \item
            \( gHg^{-1} = H \ \forall g \in G \)
    \end{enumerate}
\end{prop}

\begin{itemize}
    \item
        \( 1 \Rightarrow 2 \)

        Рассмотрим \( h \in H, g \in G \)

        \( gH = Hg \Rightarrow gh = h'g \), \( h' \in h \)

        Тогда \( g h g^{-1} = h' g g^{-1} = h' \in H \)
    \item
        \( 2 \Rightarrow 3 \)

        \( g H g^{-1} \subseteq H \)

        \( \forall h \in H \) имеем \( h = (g g^{-1}) h (g g^{-1}) = g (g^{-1} h g) g^{-1} \in g H g^{-1} \)

        А это значит, что \( H \subseteq g H g^{-1} \Rightarrow H = g H g^{-1} \)
    \item
        \( 3 \Rightarrow 1 \)

        \( g H g^{-1} = H \)

        Домножим справа на \( g \): \( gH = Hg \)
\end{itemize}

\textbf{Цель}: определить факторгруппу, то есть ввести операцию на множестве левых смежных классов \( \fc{G}{H} \),
где \( G \) --- группа, а \( H \subseteq G \) --- подгруппа.

Хочется чего-нибудь интуитивного: \( (g_1 H) (g_2 H) = g_1 g_2 H \).

Проверим необходимые свойства:
\begin{enumerate}
    \item
        Ассоциативность: \( ( (g_1 H) (g_2 H) ) (g_3 H) = g_1 g_2 g_3 H = (g_1 H) ( (g_2 H) (g_3 H) ) \)
    \item
        Нейтральный элемент: \( eH = H \quad (eH) (gH) = (gH) = (gH) (eH) \)
    \item
        Обратный элемент: \( (gH)^{-1} = g^{-1} H \quad (g^{-1} H) (g H) = eH = (g H) (g^{-1} H) \)
\end{enumerate}

Осталось самое нетривиальное --- корректность операции.
Надо как-то убедиться, что операция не изменяется при выборе разных представителей \( g_1, g_2 \) смежных классов:
\begin{gather*}
    (g_1 h_2) (g_2 h_2) H = g_1 g_2 H
    \\
    h_1 g_2 (h_2 H) = g_2 H
    \\
    g_2^{-1} h_1 g_2 H = H
    \\
    \Updownarrow
    \\
    g_2^{-1} h_1 g_2 \in H
    \\
    \Updownarrow
    \\
    \text{\( H \) нормальна}
\end{gather*}

\begin{defn}{Факторгруппа}{}
    Факторгруппа группы \( G \) по нормальной подгруппе \( H \)
    называется множество (левых) смежных классов \( \fc{G}{H} \) с операцией
    \( (g_1 H) (g_2 H) = g_1 g_2 H \)
\end{defn}

\begin{example}{}{}
    \( G = (\ZZ, +) \), \( H = n \ZZ \)

    \( \fc{G}{H} = \{ s + H = s + n \ZZ = \ol{s} \} = \ZZ_n \)
\end{example}

\begin{defn}{Гомоморфизм}{}
    Пусть \( G \) и \( F \) --- группы.
    Тогда отображение \( \varphi : G \to F \) называется гомоморфизмом,
    если \( \varphi(a \cdot b) = \varphi(a) \odot \varphi(b) \) для любых \( a, b \in G \).
\end{defn}

\begin{lemma}{}{}
    Если \( \varphi : G \to F \) --- гомоморфизм, то:
    \begin{enumerate}
        \item
            \( \varphi(e_G) = e_F \)
        \item
            \( \varphi(g^{-1}) = \varphi(g)^{-1} \) для любого \( g \in G \)
    \end{enumerate}
\end{lemma}

\begin{enumerate}
    \item
        \(
            \varphi(e_G) = \varphi(e_G \cdot e_G) = \varphi(e_G) \odot \varphi(e_G)
            \Rightarrow
            e_F = \varphi(e_G)
        \)
    \item
        \( \varphi(e_G) = \varphi(g \cdot g^{-1}) = \varphi(g) \odot \varphi(g^{-1}) \)

        \( e_F \varphi(g) \odot \varphi(g^{-1}) \Rightarrow \varphi(g^{-1}) = \varphi(g)^{-1} \)
\end{enumerate}

\begin{defn}{Биекция}{}
    Гомоморфизм \( \varphi : G \to F \) называется изоморфизмом,
    если \( \varphi \) --- биекция.
\end{defn}

\begin{exercise}{}{}
    Существует \( \varphi^{-1} : F \to F \) --- тоже гомоморфизм.
\end{exercise}

\begin{defn}{Изоморфность групп}{}
    Группы \( G \) и \( F \) изоморфны,
    если существует изоморфизм \( \varphi : G \to F \).
\end{defn}

\begin{thrm}{}{}
    \begin{enumerate}
        \item
            Всякая бесконечная циклическая группа изоморфна \( (\ZZ, +) \).
        \item
            Всякая конечная циклическая группа порядка \( n \) изоморфна \( (\ZZ_n, +) \).
    \end{enumerate}
\end{thrm}

\begin{enumerate}
    \item
        \( G = \langle g \rangle \) и \( |G| = \infty \)

        \( g^m \xmapsto{\varphi} m \)
    \item
        \( G = \langle g \rangle \) и \( |G| = n \)

        \( g^m \xmapsto{\varphi} \ol{m} \pmod n \)
\end{enumerate}

\begin{example}{}{}
    \( (\RR, +) \cong (\RR_{> 0}, \times) \)

    \( x \mapsto e^x \)
\end{example}

\begin{defn}{Ядро и образ}{}
    С каждым гомоморфизмом \( \varphi : G \to F \) связаны:
    \begin{itemize}
        \item
            \( \ker \varphi = \{ a \in G \: \vline \: \varphi(a) = e_F \} \)
        \item
            \( \Img \varphi = \{ b \in F \: \vline \: \exists a \in G : \varphi(a) = b \} \)
    \end{itemize}
\end{defn}

\begin{exercise}{}{}
    \( \ker \varphi \subseteq G \) и \( \Img \varphi \subseteq F \) --- подгруппы.
\end{exercise}

\begin{lemma}{}{}
    Гомоморфизм \( \varphi : G \to F \) инъективен \( \Leftrightarrow \ker \varphi = \{ e_G \} \)
\end{lemma}

\begin{itemize}
    \item
        \( \Rightarrow \)

        Если \( e_G \neq g \in \ker G \), то \( \varphi(g) = \varphi(e_G) = e_F \) --- элементы ``склеились''.
    \item
        \( \Leftarrow \)

        Пусть \( \varphi(g_1) = \varphi(g_2) \).

        \( \varphi(g_1 g_2^{-1}) = \varphi(g_1) \odot \varphi(g_2^{-1}) = \varphi(g_1) \odot \varphi(g_2)^{-1} = e_F \).

        Значит \( g_1 g_2^{-1} \in \ker \varphi \Rightarrow g_1 g_2^{-1} = e_G \Rightarrow g_1 = g_2 \).
\end{itemize}

\begin{coroll}{}{}
    \( \varphi \) --- биективен \( \Leftrightarrow \) \( \ker \varphi = \{ e_G \} \) и \( \Img \varphi = F \).
\end{coroll}

\begin{prop}{}{}
    \( \ker \varphi \) --- нормальная подгруппа.
\end{prop}

Проверим, что \( \forall g \in G \), \( \forall h \in \ker \varphi \) выполнено \( g h g^{-1} \in \ker \varphi \):
\[
    \varphi(g h g^{-1})
    =
    \varphi(g) \odot \varphi(h) \odot \varphi(g^{-1})
    =
    \varphi(g) \odot e_F \odot \varphi(g)^{-1}
    =
    e_F
\]

\begin{thrm}{Теорема о гомоморфизме}{}
    Для любого гомоморфизма \( \varphi : G \to F \) верно \( \Img \varphi \cong \fc{G}{\ker \varphi} \).
\end{thrm}

\begin{example}{}{}
    \( G = (\RR, +) \), \( H = (\ZZ, +) \), \( \fc{G}{H} \cong S^1 \).

    Это можно показать через такой гомоморфизм:
    \[
        \begin{aligned}
            \varphi : \RR &\to (\CC \setminus \{ 0 \}, \times)
            \\
            a &\to e^{2 \pi i a}
        \end{aligned}
    \]

    Нетрудно убедиться, что \( \ker \varphi = \ZZ \), а \( \Img \varphi = S^1 \).
\end{example}

\begin{exercise}{}{}
    \( G = (\RR^2, +) \), \( H = (\ZZ^2, +) \), \( \fc{G}{H} \cong T^2 \) (тор)
\end{exercise}

\begin{exercise}{}{}
    Пусть \( G \) --- группа, \( H \subseteq G \) --- подгруппа.

    Тогда \( g_1 H = g_2 H \Leftrightarrow g_1^{-1} g_2 \in H \).
\end{exercise}

Вернемся к доказательству теоремы.

Построим отображение \( \Psi : \fc{G}{\ker \varphi} \to \Img \varphi \) --- \( \Psi( g \ker \varphi) = \varphi(g) \)

\begin{itemize}
    \item
        Корректность

        \( \varphi (gh) = \varphi(g) \varphi(h) = \varphi(g) \)
    \item
        Сюръективность

        Есть по построению.
    \item
        Инъективность:
        \begin{gather*}
            \Psi(g_1 \ker \varphi) = \Psi(g_2 \ker \varphi)
            \\
            \Updownarrow
            \\
            \varphi(g_1) = \varphi(g_2)
            \\
            \Updownarrow
            \\
            g_1^{-1} g_2 \in \ker \varphi
            \\
            \Updownarrow
            \\
            g_1 \ker \varphi = g_2 \ker \varphi
        \end{gather*}
    \item
        Операция:
        \begin{gather*}
            \Psi( (g_1 \ker \varphi) (g_2 \ker \varphi) )
            =
            \Psi(g_1 g_2 \ker \varphi)
            =
            \varphi(g_1 g_2)
            =
            \\
            =
            \varphi(g_1) \varphi(g_2)
            =
            \Psi(g_1 \ker \phi) \Psi(g_2 \ker \varphi)
        \end{gather*}
\end{itemize}

\begin{defn}{Центр группы}{}
    \[
        Z(G) = \{ a \in G \: \vline \: ab = ba \ \forall b \in G \}
    \]
\end{defn}

\begin{example}{}{}
    \begin{itemize}
        \item
            \( Z(G) = G \Leftrightarrow G \) --- абелева.
        \item
            \( Z(S_n) = \{ e \} \) (при \( n \geq 3 \))
    \end{itemize}
\end{example}

\begin{prop}{}{}
    \( Z(G) \subseteq G \) --- нормальная подгруппа.
\end{prop}

Убедимся, что \( Z(G) \) --- группа.
Хотим \( \forall a, b \in Z(G) \Rightarrow a b^{-1} \in Z(G) \):
\[
    (a b^{-1}) c
    =
    a b^{-1} (c^{-1})^{-1}
    =
    a (c^{-1} b)^{-1}
    =
    a (b c^{-1})^{-1}
    =
    (b c^{-1})^{-1} a
    =
    c b^{-1} a
    =
    c (a b^{-1})
\]

Проверим нормальность, то есть хотим \( \forall a \in Z(G), b \in G \Rightarrow b a b^{-1} \in Z(G) \):
\[
    b a b^{-1} = a b b^{-1} = a \in Z(G)
\]

\begin{defn}{Прямое произведение}{}
    Даны группы \( G_1, \ldots, G_m \).

    Строим \( G_1 \times G_2 \times \ldots \times G_m = \{ (g_1, \ldots, g_m), \: g_i \in G_i \} \).

    \( (g_1, \ldots, g_m) (g_1', \ldots, g_m') = (g_1 \odot_1 g_1', \ldots, g_m \odot_m g_m') \)

    \( e = (e_1, \ldots, e_m) \)

    \( g^{-1} = (g_1^{-1}, \ldots, g_m^{-1}) \)
\end{defn}

\begin{remark}{}{}
    \begin{itemize}
        \item
            Прямое произведение коммутативно \( \Leftrightarrow \) все \( G_i \) коммутативны.
        \item
            \( |G_1 \times \ldots \times G_m| = |G_1| \ldots |G_m| \)
    \end{itemize}
\end{remark}

\begin{thrm}{Теорема о факторизации по сомножителям}{}
    Пусть \( H_i \subseteq G_i \) --- нормальные подгруппы.

    Тогда \( H_1 \times \ldots \times H_m \subseteq G_1 \times \ldots \times G_m \)
    --- тоже нормальная подгруппа, причем
    \[
        \fc{G_1 \times \ldots \times G_m}{H_1 \times \ldots \times H_m}
        =
        \fc{G_1}{H_1} \times \ldots \times \fc{G_m}{H_m}
    \]
\end{thrm}

Нормальность явно следует из покомпонентной нормальности.

Изоморфизм (нужно так же проверить \( 4 \) условия):
\[
    (g_1, \ldots, g_m) (H_1 \times \ldots \times H_m) \mapsto (g_1 H_1, \ldots, g_m H_m)
\]

\begin{thrm}{}{}
    Пусть \( n = ml \) и \( (m, l) = 1 \).

    Тогда \( \ZZ_n \cong \ZZ_m \times \ZZ_l \).
\end{thrm}

Рассмотрим отображение \( \varphi : \ZZ_n \to \ZZ_m \times \ZZ_l \):
\[
    k \bmod n \to (k \bmod m, k \bmod l)
\]

\begin{itemize}
    \item
        Корректность

        \( k' = k + ns = k + mls \)
    \item
        Гомоморфизм

        По определению операции
    \item
        Инъективность

        Проверим, что \( \ker \varphi = 0 \).

        Пусть \( k \bmod n \to (0 \bmod m, 0 \bmod l \Rightarrow m, l \: \vline \: k \Rightarrow k \vdots ml = n \)
    \item
        Сюръективность

        \( |\ZZ_n| = n = ml = |\ZZ_m \times \ZZ_l| \)

        Поэтому из инъективности следует биекция.
\end{itemize}

\begin{coroll}{}{}
    Пусть \( n \geq 2 \) и \( n = p_1^{\alpha_1} \cdot \ldots \cdot p_k^{\alpha_k} \) --- разложение на простые множители.

    Тогда \( \ZZ_n \cong \ZZ_{p_1^{\alpha_1}} \times \ldots \times \ZZ_{p^{\alpha_k}} \)
\end{coroll}

\subsection{Теория абелевых групп}

\textbf{Цель:}
описание конечных абелевых групп.

\begin{defn}{}{}
    Пусть \( (A, +) \) --- абелева группа. Рассмотрим \( a \in \ZZ \):
    \[
        s \cdot a = \left\{
            \begin{aligned}    
                & \underbrace{a + \ldots + a}_{s} &, s > 0
                \\
                & 0 &, s = 0
                \\
                & \underbrace{ - a - \ldots - a}_{|s|} &, s < 0
            \end{aligned}
        \right.
    \]
\end{defn}

\begin{defn}{Конечно порожденная абелева группа}{}
    Абелева группа \( A \) называется конечно порожденной,
    если
    \[
        \exists a_1, \ldots, a_n \in A : \forall a \in A \ a = s_1 a_1 + \ldots + s_n a_n, s_i \in \ZZ
    \]
\end{defn}

\begin{remark}{}{}
    \begin{itemize}
        \item
            Всякая конечно порождаемая абелева группа не более чем счетна.
        \item
            Всякая конечная абелева группа является конечно порожденной.
    \end{itemize}
\end{remark}

\begin{defn}{Свободная абелева группа}{}
    Конечно порожденная группа \( A \) свободна,
    если в ней существует базис:
    \begin{gather*}
        \exists a_1, \ldots, a_n \in A :
        \\
        \forall a \in A \ a = s_1 a_1 + \ldots + s_n a_n \ \text{(единственным образом)}
    \end{gather*}
\end{defn}

\begin{example}{}{}
    \( \ZZ_n \) конечно порожденная, не свободна.

    Можно представить \( 0 \) больше чем одним способом:
    \[
        0 = s_1 a_1 + \ldots + s_k a_k = (s_1 + n) a_1 + \ldots + (s_k + n) a_k
    \]
\end{example}

\begin{example}{Свободная группы}{}
    \( \ZZ^n = \{ (x_1, \ldots, x_n) \: \vline \: x_i \in \ZZ \} \)

    Стандартный базис:
    \begin{gather*}
        (1, 0, 0, \ldots, 0)
        \\
        (0, 1, 0, \ldots, 0)
        \\
        \vdots
        \\
        (0, 0, 0, \ldots, 1)
    \end{gather*}
\end{example}

\begin{prop}{}{}
    Всякая свободная абелева группа изоморфна \( \ZZ^n \) 
\end{prop}

Пусть \( A \) --- свободна, а \( \{ a_1, \ldots, a_n \} \) --- ее базис:
\[
    a = s_1 a_1 + \ldots + s_n a_n \longleftrightarrow (s_1, \ldots, s_n) \in \ZZ^n \text{ --- биекция}
\]

То, что это гомоморфизм --- очевидно.

\begin{defn}{Ранг}{}
    Ранг свободной абелевой группы \( \rk A \) --- число элементов в базисе.
\end{defn}

\begin{prop}{Корректность ранга}{}
    Любые два базиса свободной абелевой группы равномощны.
\end{prop}

Докажем от противного.
Пусть нашлись два базисва \( \{ a_1, \ldots, a_n \} \) и \( \{ b_1, \ldots, b_m \} \), для удобства \( m > n \).
Выразим \( b_i \) через \( a_j \):
\begin{gather*}
    b_1 = s_{11} a_1 + \ldots + s_{1n} a_n
    \\
    \vdots
    \\
    b_m = s_{m1} a_1 + \ldots + s_{mn} a_n
\end{gather*}

Рассмотрим матрицу \( S = (s_{ij}) \), \( \rk S \leq n < m \),
а значит ее строки линейно зависимы над \( \QQ \):
\[
    \exists \lambda_1, \ldots, \lambda_m \in \QQ, \text{ не все равные \( 0 \)}
    \\
    \lambda_1 (s_{11}, \ldots, s_{1m}) + \ldots + \lambda_m (s_{m1}, \ldots, s_{mn}) = 0 
\]

Можно домножить на НОК всех знаменателей \( \lambda_i \), теперь \( \lambda_i \in \ZZ \).
\begin{gather*}
    \lambda_1 b_1 + \ldots + \lambda_m b_m
    =
    \\
    =
    \lambda_1 (s_{11} + a_1 + \ldots + s_{1n} a_n)
    +
    \ldots
    +
    \lambda_m (s_{m1} + a_1 + \ldots + s_{mn} a_n)
    =
    \\
    =
    \underbrace{(\lambda_1 s_{11} + \ldots + \lambda_m s_{m1})}_{0} a_1
    +
    \ldots
    +
    \underbrace{(\lambda_1 s_{1n} + \ldots + \lambda_m s_{mn})}_{0} a_n
    =
    0
\end{gather*}

Нашли нетривиальное представление нуля --- противоречие.

Пусть \( e_1', \ldots, e_n' \) --- другой базис (любой набор векторов) в \( \ZZ^n \).
Тогда выразим его через \( e_1, \ldots, e_n \):
\[
    (e_1', \ldots, e_n') = (e_1, \ldots, e_n) \cdot C \quad C = (c_{ij}), \ c_{ij} \in \ZZ 
\]

\begin{prop}{}{}
    \( \{ e_1', \ldots, e_n' \} \) --- базис \( \Leftrightarrow \det C = \pm 1 \).
\end{prop}

\begin{itemize}
    \item
        \( \Rightarrow \)

        Пусть \( e_1', \ldots, e_n' \) --- базис, тогда:
        \[
            (e_1, \ldots, e_n) = (e_1', \ldots, e_n') D = (e_1, \ldots, e_n) CD
        \]

        \( CD = E \), а значит матрица \( C \) обратима.
        Но \( C, D \) --- целочисленные, а значит и их определитель тоже:
        \[
            \det E = \det (CD) = \det C \cdot \det D \Rightarrow \det C = \pm 1
        \]
    \item
        \( \Leftarrow \)

        Пусть \( \det C = \pm 1 \).

        Тогда по формуле для присоединенной матрицы \( C^{-1} \) целочисленная.

        Но тогда \( (e_1, \ldots e_n) = (e_1', \ldots, e_n') C^{-1} \).
        А значит все вектора выражаются через \( e_1', \ldots, e_n' \).
        Остается проверить единственность.

        Докажем от противного, пусть:
        \begin{gather*}
            a = s_1 e_1' + \ldots + s_n e_n' = s_1' e_1' + \ldots + s_n' e_n'
            \\
            \begin{pmatrix}
                e_1' & \ldots & e_n'
            \end{pmatrix}
            \begin{pmatrix}
                s_1
                \\
                \vdots
                \\
                s_n
            \end{pmatrix}
            =
            \begin{pmatrix}
                e_1' & \ldots & e_n'
            \end{pmatrix}
            \begin{pmatrix}
                s_1'
                \\
                \vdots
                \\
                s_n'
            \end{pmatrix}
            \\
            \Downarrow
            \\
            \begin{pmatrix}
                e_1 & \ldots & e_n
            \end{pmatrix}
            C
            \begin{pmatrix}
                s_1
                \\
                \vdots
                \\
                s_n
            \end{pmatrix}
            =
            \begin{pmatrix}
                e_1 & \ldots & e_n
            \end{pmatrix}
            C
            \begin{pmatrix}
                s_1'
                \\
                \vdots
                \\
                s_n'
            \end{pmatrix}
            \\
            \Downarrow
            \\
            C
            \begin{pmatrix}
                s_1
                \\
                \vdots
                \\
                s_n
            \end{pmatrix}
            =
            C
            \begin{pmatrix}
                s_1'
                \\
                \vdots
                \\
                s_n'
            \end{pmatrix}
            \\
            \Downarrow
            \\
            \begin{pmatrix}
                s_1
                \\
                \vdots
                \\
                s_n
            \end{pmatrix}
            =
            \begin{pmatrix}
                s_1'
                \\
                \vdots
                \\
                s_n'
            \end{pmatrix}
        \end{gather*}
\end{itemize}

\begin{example}{}{}
    \( \ZZ^1 \), \( e_1' = 2 \), \( e_2' = 3 \).

    Данный набор линейно зависим (\( 2 \cdot 3 - 3 \cdot 2 = 0 \)),
    но из него невозможно выбрать порождающее подмножество.
\end{example}

\begin{thrm}{}{}
    Всякая подгруппа \( N \) свободной абелевой группы \( L \) ранга \( n \)
    является свободной группой ранга \( \leq n \).
\end{thrm}

Сделаем индукцию по \( n \):
\begin{itemize}
    \item
        База \( n = 0 \):

        Для нулевой группы утверждение, очевидно, верно.
    \item
        Переход \( < n \to n \):

        Пусть \( e_1, \ldots, e_n \) --- базис в \( L \).

        Рассмотрим подгруппу \( L_1 = \langle e_1, \ldots, e_{n - 1} \rangle \)
        --- она является свободной группа с базисом \( e_1, \ldots, e_{n - 1} \).
        Пусть \( N_1 = N \cap L_1 \) --- подгруппа в \( L_1 \),
        по предположению индукции \( N_1 \) свободно с базисом \( f_1, \ldots, f_m \) и \( m \leq n - 1 \).

        Рассмотрим гомоморфизм:
        \begin{gather*}
            \varphi : N \to \ZZ=
            \\
            a = s_1 e_1 + \ldots + s_n e_n \to s_n
            \\
            \ker \varphi = N_1
        \end{gather*}

        \( \Img \subseteq \ZZ \Rightarrow \exists k \in \ZZ_{\geq 0} : \Img \varphi = k \ZZ \)

        \begin{enumerate}
            \item
                \( k = 0 \)

                \( \Img \varphi = 0 \), \( \ker \varphi = N_1 = N \) --- свободная абелева группа ранга \( \leq n - 1 \)
            \item
                \( k > 0 \)

                \( \exists f_{m + 1} : \varphi(f_{m + 1}) = k \)

                Докажем, что \( f_1, \ldots, f_m, f_{m + 1} \) --- базис в \( N \),
                тем самым докажем условие теоремы.

                \( \forall a \in \NN : \varphi(a) = ks \Rightarrow \varphi(a - s f_{m + 1}) = 0 \)

                Поэтому \( a - s f_{m + 1} \in N_1 \), а значит выражается через \( f_1, \ldots, f_m \).

                Покажем единственность, пусть
                \[
                    s_1 f_1 + \ldots + s_m f_m + s f_{m + 1} = s_1' f_1 + \ldots + s_m' f_m + s' f_{m + 1}
                \]
                Применим \( \varphi \):
                \begin{gather*}
                    \varphi(s f_{m + 1}) = \varphi(s' f_{m + 1})
                    \\
                    sk = s' k
                    \\
                    \Downarrow
                    \\
                    s = s'
                \end{gather*}

                Значит последнее слагаемое можно сократить.
                Но тогда получили выражение некоторого вектора двумя способами через \( f_1, \ldots, f_m \) --- противоречие.
        \end{enumerate}
\end{itemize}

\begin{remark}{Отличие от линейной алгебры}{}
    Знаем, что если \( U \subsetneq V \), то \( \dim U < \dim V \)

    У нас не так: \( 2 \ZZ \subseteq \ZZ \), но \( \rk 2 \ZZ = 1 = \rk \ZZ \)
\end{remark}

\textbf{Локальная цель:} описать все подгруппы в \( \ZZ^n \).

\( n = 1 \): все подгруппы в \( \ZZ \) --- это \( k \ZZ, k \in \ZZ_{\geq 0} \)

\textit{Гипотеза:}
надо зафиксировать \( k_1, \ldots, k_n \in \ZZ_{\geq 0} \),
\( k_1 \ZZ \times \ldots \times k_n \ZZ \subseteq \ZZ^n \).

\textit{Нет:} \( \{ (a, a), a \in \ZZ \} \subseteq \ZZ \), откуда \( k_1 = k_2 = 1 \), но это не \( \ZZ \times \ZZ \).

\begin{defn}{}{}
    Для абелевых групп прямое произведение принято обозначать \( \oplus \), а не \( \times \).
\end{defn}

\begin{thrm}{Теорема о согласованных базисах}{}
    Пусть \( L \) --- свободная абелева группа, \( N \subseteq L \) --- подгруппа.
    Тогда существует базис \( e_1, \ldots e_n \) в \( L \),
    \( \exists m \leq n \) и натуральные числа \( u_1, \ldots, u_m : u_i \: \vline \: u_{i + 1} \)
    и \( u_1 e_1, \ldots, u_m e_m \) --- базис в \( N \)
    (это означает, что в базисе \( e_1, \ldots, e_n \) подгруппа \( N \) будет иметь вид
    \( u_1 \ZZ_1 \oplus \ldots \oplus u_m \ZZ \oplus \{ 0 \} \oplus \ldots \oplus \{ 0 \} \)). 
\end{thrm}

\begin{defn}{Целочисленные элементарные преобразования строк целочисленной матрицы}{}
    \begin{enumerate}
        \item
            Прибавление к одной строки матрицы другой,
            умноженной на число.
        \item
            Перестановка двух строк.
        \item
            Умножение строки на \( \pm 1 \).
    \end{enumerate}
\end{defn}

\begin{prop}{}{}
    Всякую целочисленную матрицу \( C = (c_{ij}) \) элементарными преобразованиями
    строк и столбцов можно привести к диагональном виду.
\end{prop}

Если \( C = 0 \), то все ясно.
Иначе переставим строки и столбцы так, чтобы \( c_{11} \neq 0 \).
Из-за умножения на \( -1 \) можно считать, что \( c_{11} > 0 \).

Хочу: \( c_{11} \) делит все \( c_{ij} \).
Если \( c_{1i} = c_{11} q + r \), и \( 0 < r < c_{11} \),
то сделаем \( c_{1i} = r \) вычитанием и поставим на место \( c_{11} \).
Поскольку уменьшаем \( c_{11} \) каждый раз, процесс когда-то кончится.
Можем занулить в первой строке все кроме \( c_{11} \).
Аналогично сделаем со столбцом.

Теперь, если же есть \( c_{ij} \), которые не делятся на \( c_{11} \),
то прибавим \( j \)-й столбец к первому, теперь там есть элемент,
не кратный \( c_{11} \).
Повторим вышеописанный алгоритм.
Поскольку каждый раз \( c_{11} \) уменьшается,
алгоритм когда-то кончится.

Теперь в матрице все числа делятся на \( c_{11} \),
занулим первую строку и столбец,
рекурсивно запустимся на оставшейся подматрице.

\textbf{Доказательство теоремы о согласованных базисах:}

Пусть \( e_1', \ldots, e_n' \) --- базис в \( L \),
и \( f_1', \ldots, f_m' \) --- базис в \( N \), \( m \leq n \).
Будем менять базисы:
\[
    (f_1', \ldots, f_m') = (e_1', \ldots, e_n') C
\]

Производим элементарные преобразования над \( C \).
Элементарные преобразования строк \( C \) ---
это элементарные замены базиса в \( L \),
а элементарные преобразования столбцов \( C \)
--- элементарные замены базиса в \( N \).

Приведем \( C \) к диагональному виду,
получатся нужные нам базисы.

\begin{remark}{}{}
    Числа \( u_1, \ldots, u_m \) называются
    инвариантными множителями либо матрицы \( C \), либо пары \( (L, N) \).
    Эти числа определены однозначно.
\end{remark}

\begin{coroll}{}{}
    \[
        \fc{L}{N} \cong \ZZ_{u_1} \oplus \ldots \oplus \ZZ_{u_m} \oplus \underbrace{\ZZ \oplus \ldots \oplus \ZZ}_{n - m}
    \]
\end{coroll}

Зафиксируем согласованные базисы \( e_1, \ldots, e_n \) и \( u_1 e_1, \ldots, u_m e_m \).
Тогда:
\begin{gather*}
    L = \ZZ^n, a = s_1 e_1 + \ldots + s_n e_n \longleftrightarrow (s_1, \ldots, s_n)
    \\
    N = \ZZ^m, b = c_1 u_1 e_1 + \ldots + c_m u_m e_m
\end{gather*}
\begin{gather*}
    \fc{L}{N}
    =
    \fc{\ZZ^n}{\ZZ^m}
    =
    \fc{ \ZZ \oplus \ldots \oplus \ZZ }{ u_1 \ZZ \oplus \ldots \oplus u_m \ZZ \oplus \{ 0 \} \oplus \ldots \oplus \{ 0 \} }
    \cong
    \\
    \cong
    \fc{\ZZ}{u_1 \ZZ} \oplus \ldots \oplus \fc{\ZZ}{u_m \ZZ} \oplus \fc{\ZZ}{ \{ 0 \} } \oplus \ldots \oplus \fc{\ZZ}{ \{ 0 \} }
\end{gather*}

\begin{thrm}{}{}
    Всякая конечно порожденная абелева группа изоморфна
    \[
        \ZZ_{p_1^{k_1}} \oplus \ldots \oplus \ZZ_{p_s^{k_s}} \oplus \ZZ \oplus \ldots \oplus \ZZ
    \]
    где \( p_i \) --- простые, возможно совпадающие числа,
    \( k_1, \ldots, k_s \) --- натуральные числа.
    Кроме того, число слагаемых \( \ZZ \), а также числа \( p_i, k_i \)
    определены по группе \( A \) однозначно.
\end{thrm}

Пусть \( a_1, \ldots, a_n \) --- система порождающих в \( A \).
Рассмотрим гомоморфизм:
\[
    \varphi : \ZZ^n \to A, (s_1, \ldots, s_n) \to s_1 a_1 + \ldots + s_n a_n
\]

Он сюръективен, так как \( a_1, \ldots, a_n \) --- порождающие.
По теореме о гомоморфизме:
\[
    \Img \varphi \cong \fc{\ZZ^n}{\ker \varphi}
    =
    \ZZ_{u_1} \oplus \ldots \oplus \ZZ_{u_m} \oplus \ZZ \oplus \ldots \oplus \ZZ
\]

Остается разложить \( u_i \) на простые.

Единственность: дополнительный материал.

\begin{coroll}{}{}
    Любая конечная абелева группа изоморфна:
    \[
        A \cong \ZZ_{p_1^{k_1}} \oplus \ldots \oplus \ZZ_{p_s^{k_s}} \cong \ZZ_{u_1} \oplus \ldots \ZZ_{u_m}
    \]
\end{coroll}

Любая конечная группа, в частности, конечно порождена.
Значит она имеет вышеописанное представление.
Но слагаемых \( \ZZ \) быть не может, так как \( \ZZ \) бесконечно.

\begin{example}{}{}
    \[
        \ZZ_6 \oplus \ZZ_{20}
        \cong
        \ZZ_2 \oplus \ZZ_3 \oplus \ZZ_{2^2} \oplus \ZZ_5
        \cong
        \ZZ_2 \oplus \ZZ_{2^2 \cdot 3 \cdot 5}
    \]
\end{example}

\begin{defn}{Экспонента}{}
    Экспонентой конечной абелевой группы \( A \)
    называется
    \[
        \exp (A) = \text{НОК}( \ord(a), a \in A )
    \]
\end{defn}

\begin{exercise}{}{}
    \[
        \exp A = \min \{ m \in \NN \: \vline \: ma = 0 \ \forall a \in A \}
    \]
\end{exercise}

\begin{prop}{}{}
    \( \exp (A) = u_m \) --- последний инвариантный множитель.
\end{prop}

\( A = \ZZ_{u_1} \oplus \ldots \oplus \ZZ_{u_m}, u_i \: \vline \: u_{i + 1} \Rightarrow u_m a = 0 \ \forall a \in A \)

С другой стороны в \( (\ol{0}, \ldots, \ol{0}, \ol{1}) \in A \),
и его порядок ровно \( u_m \)

\begin{coroll}{}{}
    \( \exp (A) \) равно наибольшему порядку элемента.
\end{coroll}

\begin{coroll}{Критерий цикличности}{}
    Группа \( A \) циклическая \( \Leftrightarrow \exp(A) = |A| \)
\end{coroll}

\subsection{Действия групп}

\begin{defn}{Действие группы}{}
    Действие группы \( G \) на множестве \( X \)
    это отображение \( G \times X \to X : (g, x) \to gx \),
    такое что
    \begin{enumerate}
        \item
            \( ex = x \ \forall x \in X \)
        \item
            \( g (hx) = (gh) x \)
    \end{enumerate}

    % \( G \) действует на \( X \): \( G \hookdownarrow X \)
\end{defn}

Задание действия определяет биекции \( \alpha_g : X \to x, x \to gx \)

Пусть \( S(x) \) --- группа всех биекций на множестве \( X \).
Например, если \( |X| = n \), то \( S(X) \cong S_n \).

Действие \( G \) на \( X \) \( \Leftrightarrow \) гомоморфизм \( \alpha : G \to S(x) \)
\begin{itemize}
    \item
        \( \Rightarrow \)

        \( g \to \alpha_g \)

        \( \alpha_{gh} (x) = (gh) (x) = g(hx) = \alpha_g ( \alpha_h x ) \)
    \item
        \( \Leftarrow \)

        \( \alpha : G \to S(x) \Rightarrow \):
        \[
            \begin{aligned}
                G \times X &\to X
                \\
                (g, x) &\mapsto \alpha_g (x)
            \end{aligned}
        \]

        (Выполнены аксиомы действия)
\end{itemize}

\begin{example}{}{}
    \( G = S_n \)

    \( X = \{ 1, \ldots, n \} \)

    \( (\sigma, i) \to \sigma(i) \)
\end{example}

\begin{defn}{Орбита}{}
    Орбита точки \( x \in X \) под действием группы \( G \) это
    \[
        Gx = \{ y \in X \: \vline \: \exists g \in G : gx = y \}
    \]
\end{defn}

\begin{remark}{}{}
    ``\( x \) и \( x' \) лежат в одной \( G \)-орбите''
    является отношением эквивалетности.
\end{remark}

\begin{defn}{}{}
    Стабилизатор (стационарная подгруппа) точки \( x \in X \) это
    \[
        St(x) = \{ g \in G : gx = x \}
    \]
\end{defn}

\begin{example}{}{}
    \( G = S^1 \), \( X = \CC \)

    \( (z, w) \to zw \)

    \(
        St(w) = \begin{cases}
            \begin{aligned}
                \{ 1 \}, & w \neq 0
                \\
                G,       & w = 0
            \end{aligned}
        \end{cases}
    \)
\end{example}

\begin{example}{}{}
    \( G = SL_n (\RR) \), \( X = \RR^n \), \( (A, v) \to Av \)

    Сколько орбит?

    \( n = 1 \Rightarrow \inf \) орбит

    Иначе орбиты \( \{ 0 \} \) и \( \RR \setminus 0 \)
\end{example}

\begin{lemma}{}{}
    \( G \) конечна.
    Тогда \( \forall x \in X \ |Gx| = \frac{|G|}{St(x)} \).
\end{lemma}

Нетрудно убедиться, что есть такая биекция:
\[
    \begin{aligned}
        \fc{G}{St(x)} &\to Gx
        \\
        g St(x) &\mapsto gx
    \end{aligned}
\]

\begin{enumerate}
    \item
        Корректность

        \( \forall h \in St(x) \ ghSt(x) \mapsto (gh) x = g (hx) = gx \)
    \item
        Сюръективность
    \item
        Инъективность

        \(
            g x = g' x
            \Rightarrow
            x = g^{-1} g' x
            \Rightarrow
            g^{-1} g' \in St(x)
            \Rightarrow
            g St(x) = g' St(x)
        \)
\end{enumerate}

\begin{defn}{}{}
    \begin{enumerate}
        \item
            Действие \( G \) на \( X \) транзитивно,
            если есть всего одна орбита
            \( \Leftrightarrow \exists g \in G : gx = x' \)
        \item
            Действие \( G \) на \( X \) свободно,
            если \( gx = x \) для некоторого \( x \in X \)
            влечет \( g = e \) \( \Leftrightarrow \forall x \in X \ St(x) = \{ e \} \)
        \item
            Действие \( G \) на \( X \) эффективно,
            если условие \( gx = x \ \forall x \in X \) влечет \( g = e \)
            \( \Leftrightarrow \bigcap\limits_{x \in X} St(x) = \{ e \} \)

            Свободность \( \Rightarrow \) эффективность.
    \end{enumerate}
\end{defn}

\begin{example}{}{}
    \( S_n \) на \( \{ 1, \ldots, n \} \) транзитивно, несвободно (при \( n \geq 3 \)) и эффективно
\end{example}

\begin{example}{}{}
    \( S^1 \) на \( \CC \) нетранзитивно, несвободно и эффективно.
\end{example}

\begin{example}{}{}
    \( SL_n(\RR) \) на \( \RR^n \) нетранзитивно, несвободно и эффективно.
\end{example}

\begin{defn}{}{}
    Ядро неэффективности \( K = \bigcap\limits_{x \in X} St(x) \)
\end{defn}

\( \alpha G \to S(x) \Rightarrow K = \ker \alpha \Rightarrow K \) нормально \( \Rightarrow \exists \) факторгруппа \( \fc{G}{K} \)
